\documentclass[11pt]{article}
\usepackage[margin=1in]{geometry}
\usepackage{booktabs}
\usepackage{amsmath}
\setlength{\parskip}{4pt}
\title{Control-Flow Work on an RV32I Core\\\large From a plain 5-stage pipeline to a 21\% faster one}
\author{CTRL-ALT-ELITE}
\date{November 2026}
\begin{document}
\maketitle

\begin{abstract}
This report covers five experiments run on a five-stage in-order RV32I core: a fetch-stage
pre-decoder, a branch history table and branch target buffer, a repair to that table after it
turned out to be hurting performance, macro-op fusion, and early branch resolution. Two of the
five were kept. Fusion was built, verified, measured on the graded workload, found not to
transfer, and dropped; a second variant was built to test the obvious explanation for that
failure and it did not transfer either. The two that were kept take the graded benchmark suite
from 1.950\,s to 1.546\,s, a 20.7\% reduction, with all eleven programs still correct. Along
the way the project grew a small assembler, a fourteen-kernel benchmark suite, and three
analysis tools, because none of the questions being asked could be answered with the tests that
existed.
\end{abstract}

\section{Starting point}

The core at the start of this work was a five-stage in-order pipeline (IF, D, X, M, WB) with
a combinational instruction memory, no branch prediction of any kind, and 53 directed tests
passing. Every conditional branch and every jump waited until execute to find out where it was
going. When a branch mispredicted, the two instructions behind it in F and D had to be
discarded and the pipeline refilled: two wasted cycles per redirect.

Everything in this report is about control flow, because that is where the cycles were. The
arithmetic path was not the problem; the pipeline was spending its time waiting on decisions
that could have been made earlier.

Two constraints shaped the work. The graded benchmark programs are compiled binaries that the
course runs on its own autograder. Their sources are not available, so every claim has to be
checked twice: once on code we can see, and once on the submission that gets graded. The second
check is the one that decides whether a feature ships.

\section{How everything was measured}

\subsection{The experiment loop}

Each change went through the same steps. A baseline was frozen: the full suite run once, with
pass/fail and per-test cycles recorded, so that any later run could be compared against a
known-good state. Then a hypothesis was written down before any RTL, stating what should change
and what must not. Then the module's ports and timing contract were specified. Then the RTL was
written, unit-tested, and put through the full suite. Any test that passed at the baseline and
fails afterwards is a behavioural divergence, and a cycle count that moves on a test the change
should not touch is treated the same way.

\subsection{Feature masks}

Two of the features in this report can be switched off at compile time, by a single
\texttt{localparam}:

\begin{verbatim}
    localparam bit FUSE_ENABLE = 1'b0;   // macro-op fusion
\end{verbatim}

That one bit made the whole difference in the fusion work. With it low, the design is
bit-identical to the core without the feature: the same 67 tests pass with the same cycle
counts, test by test. That let the plumbing be landed, and proven inert, before the feature was
allowed to change behaviour. It also made every later comparison a single-character diff against
the same build, which removes a class of ambiguity that otherwise takes hours to chase.

\subsection{The autograder as the second measurement}

Local measurements use 14 benchmark kernels written for this project. They are useful and they
are not the graded workload. Every feature that looked promising locally was therefore also
built into a submission zip and graded. Twice this second step disagreed with the first, in
opposite directions, and both times it was right.

\subsection{What a number means}

Simulation is deterministic. The same RTL and the same program image produce the same cycle
count, always. So a cycle difference between two builds is attributable, and there is no
run-to-run noise to average away.

The autograder is not like that. It reports simulation wall time rounded to a millisecond. On
a program that runs for 12\,ms, that is an 8\% quantum, so small programs cannot resolve small
effects at all. The programs worth reading are the ones above 50\,ms: Perf, Fibonacci, Bubble
sort, Check vowel, Factorial, GCD and Sum array.

\section{Tools and benchmarks built along the way}

None of the questions in this report could be answered with the tests that existed, so a fair
amount of the work was building the means to measure.

\subsection{An assembler}

\texttt{tools/rvasm.py} is a small RV32I assembler. The need for it was practical: the existing
directed tests were hand-encoded, and hand encoding is a real source of error. During this work
a helper that builds \texttt{jal} placed \texttt{imm[11]} at instruction bit 10 instead of bit
20, corrupting the destination register field as well. Reading the source would not have found
that.

The assembler is validated by reproduction: it assembles all thirteen hand-written sources in
\texttt{pipeline\_tests} and reproduces their checked-in images byte for byte. That
covers the instruction subset those tests use, and any error in a field would show as a diff.

\subsection{Fourteen kernels}

\texttt{code/benchmarks/} holds programs whose sources sit next to their images. They were
written for specific measurements:

\begin{itemize}
\item \texttt{sweep4} through \texttt{sweep100}: one counted loop, one branch in the body, five
  different trip counts. The only variable between them is iterations.
\item \texttt{loop10}, \texttt{fib}: loops with arithmetic in the body, for forwarding and
  dependency work.
\item \texttt{nested}: five outer iterations of a four-iteration inner loop. The inner loop
  exits once per outer iteration, which no single loop can expose.
\item \texttt{if-in-loop}: a branch inside a hot loop that falls through exactly once.
\item \texttt{sum-array}, \texttt{array-idx}: array loops with load-use pairs and, in the
  first, stores.
\item \texttt{smul}: a shift-and-add multiply over 32 iterations, which is the shape a libgcc
  multiply takes on this ISA.
\item \texttt{callret}: \texttt{jal x1} and \texttt{jalr x0,0(x1)} pairs, the canonical
  call/return idiom for a return-address stack.
\item \texttt{smoke}: two instructions, to separate harness failures from program failures.
\end{itemize}

\subsection{Analysis tools}

Three scripts do the counting that the design decisions depended on:

\begin{itemize}
\item \texttt{fusion\_analysis.py} recovers the dynamic instruction stream from each kernel's
  source plus its measured retired-instruction count, so a pair inside a loop is counted once
  per iteration rather than once per line.
\item \texttt{objdump\_analysis.py} reads a RISC-V \texttt{objdump} listing and reports the
  instruction mix and every \texttt{addi}+branch adjacency, split by which operand the branch
  tests, with loop-backs called out separately. It exists because the hand-written kernels
  turned out not to be representative, and compiled code had to be inspected directly.
\item The commit-trace comparison logs every architectural writeback from two builds and diffs
  the sequences. It is the tool that found all three bring-up bugs in the early-resolution work.
\end{itemize}

\subsection{Unit testbenches}

Two modules have their own testbenches, each a table of input rows with expected outputs, driven
on one clock edge and checked on the other with case equality:

\begin{itemize}
\item \texttt{pre\_decode\_tb}: 27 rows, all passing.
\item \texttt{fuse\_detect\_tb}: 12 rows for v1 of the fusion recogniser, 16 after v2 added the
  mirrored operand form. Includes a second base address and a row whose source and destination
  registers differ, so that neither the PC nor the register fields can go unread without a
  failure.
\end{itemize}

\section{Experiment 1: a fetch-stage pre-decoder}

\subsection{Hypothesis}

Every jump waited for execute. But the instruction word is available in IF, and the instruction
memory is combinational, so a \texttt{jal} can be recognised and its target computed in the same
cycle it is fetched. The prediction was that \texttt{jal} would lose four cycles and everything
else would not move.

\subsection{Build}

\texttt{pre\_decode.sv} takes the instruction, its PC and a validity signal, and produces the
target, the link value, a CFI class, and three qualifiers: is this a call, is it a return, does
it redirect from IF. It is purely combinational. The redirect itself is one term in the
fetch-PC multiplexer.

\subsection{Result}

\begin{center}
\begin{tabular}{lrrr}
\toprule
test & baseline & with pre-decoder & change \\
\midrule
\texttt{rv32ui-p-jal}   & 27 & 23 & $-4$ \\
\texttt{rv32ui-p-auipc} & 30 & 26 & $-4$ \\
\texttt{test1}          & 53 & 47 & $-6$ \\
\texttt{test2}          & 64 & 56 & $-8$ \\
\texttt{rv32ui-p-jalr}  & 98 & 98 & $0$ (predicted exactly zero) \\
\bottomrule
\end{tabular}
\end{center}

The other 49 tests were unchanged and all 53 passed. The \texttt{jalr} row is the useful one:
its target depends on a register value, so IF cannot compute it, and the hypothesis said it
would not move. It did not.

\section{Experiment 2: a branch history table and a branch target buffer}

\subsection{Hypothesis and build}

Conditional branches are the remaining redirects. A two-bit counter per branch, indexed by PC,
can predict direction; a small table of stored targets can predict \texttt{jalr} targets. Both
are read in IF and trained in execute.

\subsection{Result}

The two halves behaved very differently.

\begin{center}
\begin{tabular}{lrrrr}
\toprule
test & cycles, before & cycles, after & flushes, after & verdict \\
\midrule
\texttt{rv32ui-p-jalr} & 98 & \textbf{92} & 12 & BTB pays: $-6$ cycles, $-3$ flushes \\
\texttt{rv32ui-p-bne}  & 315 & 319 & 31 & worse \\
\texttt{rv32ui-p-beq}  & 311 & 317 & 30 & worse \\
\texttt{rv32ui-p-blt}  & 311 & 317 & 30 & worse \\
\texttt{rv32ui-p-bge}  & 347 & 357 & 41 & worse \\
\texttt{rv32ui-p-bgeu} & 372 & 382 & 41 & worse \\
\bottomrule
\end{tabular}
\end{center}

The \texttt{jalr} improvement is what a working BTB should look like. The branch tests getting
\emph{slower} is not what a working direction predictor looks like. The BTB was adopted; the BHT
was not mergeable as it stood.

\section{Experiment 3: why the branch table was hurting}

\subsection{The suspect}

With 64 untagged entries indexed by \texttt{pc[7:2]}, each entry covers 256 bytes of program.
A test with many branches in a small region will have several of them mapping to the same
counter. The obvious fix was more entries.

The index was raised from 6 bits to 10. Nothing changed: the branch tests stayed exactly as
slow. That result is worth recording because the first attempt at it silently did nothing, since
the instantiation in \texttt{rv\_core.sv} overrode the module's parameter. Once the override was
changed, the numbers were real: 1024 entries, same regression. So aliasing was not the
mechanism, or not the whole mechanism.

\subsection{The actual cause}

The table had a tag per entry, and the read path compared the tag and a valid bit. That was
working. The problem was on the write side: the tag recorded which branch had written the entry
last, but the two-bit counter held a mixture of outcomes from whichever branches had used that
entry before. A read whose tag matched therefore got a counter that did not describe it.

The fix is a write-side tag comparison as well, plus cold seeding: a new entry starts at
\texttt{STRONG\_NT} and takes only the current outcome, rather than inheriting whatever the
previous occupant left behind. An intermediate attempt seeded new entries as taken, which
invented predictions that did not exist and made the branch tests worse still (335 cycles, 39
flushes) before the cold seed put them back at the pre-predictor numbers (311 cycles, 27
flushes).

The finished table costs 704 flip-flops.

\subsection{What this left unproven}

The BHT made zero correct predictions on the directed tests, because those tests execute each
branch about once. The table was now correct and not harmful, but nothing in the suite could
show it helping. That is the direct origin of the benchmark suite: a predictor needs a branch
that repeats, and the directed tests had none.

\section{Experiment 4: macro-op fusion}

\subsection{Where the idea came from}

After the predictor, the remaining overhead on the kernels was 13.6\% of cycles, and the largest
single item was the redirect cost itself. The other axis is to eliminate instructions: if two
architecturally separate instructions always execute back to back, they can occupy one slot.

The obvious candidate is the loop counter. Every counted loop ends with the same two
instructions:

\begin{verbatim}
    addi x5, x5, -1
    bne  x5, x0, loop
\end{verbatim}

\subsection{The gate, before any RTL}

The first version of this analysis counted adjacencies once per source line, which overstates
the benefit: a pair inside a loop fires once per iteration. Re-running it against the recovered
dynamic stream gave the numbers that decided whether to build anything at all.

\begin{center}
\begin{tabular}{lrrr}
\toprule
candidate & dynamic pairs & share of cycles & implementation cost \\
\midrule
\texttt{addi} + conditional branch & 310 & 23.2\% & 3 of 5 \\
\texttt{lw} + consumer             & 12 & 0.9\% & 5 of 5 \\
\texttt{add} + \texttt{sw}          & 12 & 0.9\% & 2 of 5 \\
\texttt{add} + \texttt{lw}          & 12 & 0.9\% & 2 of 5 \\
\texttt{lui}/\texttt{auipc} + op    & 2 & 0.1\% & 2 of 5 \\
\texttt{slt} + branch               & 0 & 0\% & 2 of 5 \\
\texttt{lw} + \texttt{sw}           & 0 & 0\% & 3 of 5 \\
\bottomrule
\end{tabular}
\end{center}

One candidate was worth two orders of magnitude more than the others, and it was also the
cheapest to build, because the adder and comparator it needs already existed in execute.
Everything memory-side occurs only in indexed array code and only once per iteration.

The rule was then narrowed to what would actually be built. A fused pair can reference at most
two register-file reads, because the file has two ports. Requiring the branch's \texttt{rs1} to
be the addi's destination means the branch's other operand is the addi's result, not a third
read, so two reads suffice. Under that restriction, 287 of the 310 pairs survive, 25.6\% of
dynamic instructions, and 97\% of what survives is a loop-back. The predicted bound was 21.4\%
of the suite's cycles.

\subsection{Contract}

The recogniser was specified before it was written: ten ports, a four-sentence behaviour
statement, an eight-row rejection table, a timing contract, and the payload that has to ride
with the instruction once the second word is consumed. A few decisions were made deliberately
and written down:

\begin{itemize}
\item Invalid branch funct3 values (\texttt{010}, \texttt{011}) are rejected rather than fused
  with a not-taken comparator. Fusing them would be harmless, but ``fused implies a real
  branch'' is a one-line invariant for the testbench.
\item An \texttt{addi} whose destination is \texttt{x0} is rejected. Correct to fuse, pointless:
  nothing is written and the compare would be against zero.
\item The branch target is pre-computed in fetch, and \texttt{br\_imm} is carried anyway so a
  later variant can compute it in execute without changing the interface. An unused output that
  exists for a variant is not dead code.
\item Both payload outputs are always driven, even when a given variant uses only one, so the
  same testbench passes unmodified on every variant.
\end{itemize}

\subsection{v1, and what integration actually cost}

The testbench came before the RTL. Twelve rows: seven rejections from the truth table and five
acceptances, including a second base address and a pair whose source and destination registers
differ, so that neither the PC nor the register fields could go unread without a failure.
Against a stub with all outputs tied low, five rows failed and seven passed. A testbench that
cannot fail is not a testbench.

The recogniser itself was then about forty lines: field extraction, seven comparisons, a
priority encoder for the payload. It passed 16 of 16. The integration is where the work was.

Three defects, all in the wiring between stages:

\begin{enumerate}
\item \textbf{The payload stopped at decode.} The D/X register had no ports for the fused flag,
  branch condition or target, so three nets were implicitly declared one bit wide. That
  compiles. It is worse than an error: a truncated target redirects to address zero, and the
  undriven fused flag reads as zero, which quietly disables the feature and hides the other two.
\item \textbf{Execute had no fused term.} The check that compares the prediction against the
  real next PC was written as ``is a branch or is a jump''. For a fused pair the opcode is the
  addi, and neither is true, so a taken fused branch computed the right target and then never
  acted on it.
\item \textbf{The operand index was wrong from decode onward.} For a fused pair the second
  operand is the branch's \texttt{rs2}, but the pipeline carried the addi's \texttt{rs2} field,
  which for an I-type is part of the immediate. Forwarding, the load-use hazard and the
  register-file read all used that wrong index. This one failed exactly one directed test,
  \texttt{rv32ui-p-jal}, and took a commit trace to find: the branch compared a stale register
  value, took the wrong direction, and redirected into the failure path.
\end{enumerate}

\subsection{Local result}

With the feature enabled and everything else unchanged, 67 of 67 tests passed and the fourteen
kernels dropped from 1339 to 1048 cycles.

\begin{center}
\begin{tabular}{lrrrr}
\toprule
kernel & cycles, off & cycles, on & saved & flush cycles \\
\midrule
\texttt{sweep100} & 317 & 217 & $-100$ & 3 / 3 \\
\texttt{sweep50}  & 167 & 117 & $-50$ & 3 / 3 \\
\texttt{nested}   & 112 &  87 & $-25$ & 10 / 10 \\
\texttt{smul}     & 218 & 186 & $-32$ & 5 / 5 \\
\texttt{fib}      &  69 &  59 & $-10$ & 3 / 3 \\
\textbf{total}    & \textbf{1339} & \textbf{1048} & $\mathbf{-291\ (-21.7\%)}$ & unchanged \\
\bottomrule
\end{tabular}
\end{center}

The saving is exactly one cycle per fused pair on every kernel, and the flush counts are
identical test by test. That second observation is the check that mattered: fusion removes
branches from the stream the predictor is measured on, so a badly wired prediction path would
have shown up as extra flush cycles eating the gain. It did not.

\subsection{Graded result: no transfer}

The submission was graded. Every program was correct, and the total was 1.719\,s against
1.713\,s without the feature: a 0.4\% regression, indistinguishable from noise.

The measurement is sensitive enough to see what it was testing for. Prediction removes about
three cycles per loop iteration and shows as 12 to 18\% on the loop-heavy programs. Fusion
removes one instruction per pair and should have been visible at half that. It was not visible
at all, so the pairs were not being fused.

\subsection{The obvious explanation, and its refutation}

The rule required the branch's \texttt{rs1} to be the addi's destination. Compiled loops often
compare the other way round:

\begin{verbatim}
    addi x5, x5, -1
    blt  x5, x6, loop     rs1 == rd   ->  v1 fuses
    bge  x6, x5, loop     rs2 == rd   ->  v1 rejects
\end{verbatim}

Both are correct RISC-V and both appear in compiled code. The project's own kernels are all
written in the first form, which is why they measure 21.7\%: they were shaped to fit the rule
being measured. A second variant accepted both forms, with one extra bit of payload saying
which side of the comparison the addi's result belongs on. It was a strict superset: every v1
row of the testbench still passed.

It measured 1.722\,s, against 1.713\,s for prediction alone and 1.719\,s for v1. No
improvement, and Perf was 3\% slower.

Between them, the two variants cover every adjacent \texttt{addi}+branch pair in which the
branch tests the addi's result, in either operand position. The union still shows nothing. So
the operand order was not the explanation, and neither variant belongs in the shipped design.

Two candidate reasons remain and the measurements do not separate them: either the pattern is
genuinely rare in compiled loops, or the grader's metric blunts it because the fusion build adds
a second combinational instruction-memory port, which costs simulator time per cycle regardless
of architectural cycles. The two heaviest loop programs are unchanged to the millisecond, which
argues the cycle effect itself is small.

\section{Experiment 5: early branch resolution}

\subsection{The budget that pointed here}

After the predictor and the rejection of fusion, the cycle budget of the kernels looked like
this:

\begin{center}
\begin{tabular}{lr}
\toprule
category & share of cycles \\
\midrule
useful instructions & 86.4\% \\
mispredict flushes (two cycles each) & 8.3\% \\
pipeline fill and alignment & 4.2\% \\
load-use stalls & 1.2\% \\
\bottomrule
\end{tabular}
\end{center}

Fill is irreducible in a five-stage pipeline. The mispredict cost was the largest remaining
item, and unlike the predictor's accuracy it is a property of \emph{where} the decision is made.
A redirect resolved in execute discards two instructions; resolved one stage earlier, in decode,
it discards one.

\subsection{Design}

Conditional branches compare in D:

\begin{verbatim}
    d_npc        = d_brtaken ? d_pc + d_imm : d_pc + 4
    mispredict_d = (d_opcode == B_TYPE) && !hazard && (fd_pred_npc != d_npc)
\end{verbatim}

Two parts of that are not obvious. The operands are still in flight when the branch reaches D, so
D forwards from X and M. And the flush signal cannot simply be reused: today one signal clears
both the F/D and D/X registers, which is correct when detection happens in execute because
everything above the branch is younger. Detected in decode, the instruction in execute is
\emph{older} and must survive. The new signal clears the F/D register only.

\subsection{Three defects}

\begin{enumerate}
\item \textbf{A load in M.} The forwarding value for the memory stage has no memory-data case:
  it selects immediate, PC+4 or ALU result. That is exactly why the existing M$\rightarrow$X
  forwarding masks loads out and lets them reach writeback first. Resolving in D has no such
  option, so a load's data is taken from the memory output in M. This one defect accounted for
  seven of eight failing tests: every sub-word load and store.
\item \textbf{\texttt{lui}, \texttt{auipc} and \texttt{jal} in X.} The execute-stage result is
  the ALU output, but those instructions write the immediate or PC+4, selected only at
  writeback. A branch comparing against a \texttt{lui} result read zero, which is what broke
  \texttt{rv32ui-p-bgeu}: it compares \texttt{0x7fffffff} against a \texttt{lui}-built
  \texttt{0x80000000}.
\item \textbf{The execute-stage check had to be removed.} Leaving the branch opcode in the
  execute-stage prediction check makes the X redirect win over the D redirect for the same
  branch. The design is then correct, all 67 tests pass, and exactly zero cycles are saved.
  This defect was introduced during bisection and cost an hour before the trace showed the
  redirect still originating in execute.
\end{enumerate}

\subsection{Result}

\begin{center}
\begin{tabular}{lrrr}
\toprule
kernel & baseline & early resolution & saved \\
\midrule
\texttt{sweep100} & 317 & 314 & $-3$ \\
\texttt{nested}   & 112 & 102 & $-10$ \\
\texttt{if-in-loop}&  56 &  50 & $-6$ \\
\texttt{smul}     & 218 & 213 & $-5$ \\
\texttt{fib}      &  69 &  66 & $-3$ \\
\textbf{total}    & \textbf{1339} & \textbf{1288} & $\mathbf{-51\ (-3.8\%)}$ \\
\bottomrule
\end{tabular}
\end{center}

Exactly one cycle per mispredict: \texttt{nested} saves ten for its ten, the loop kernels save
three each, which is entry, exit and the cold first iteration. 67 of 67 pass.

On the graded suite the same build returned 1.677\,s against 1.713\,s for prediction alone,
$-2.1\%$, with eight programs faster, three unchanged and none slower. The three that did not
move are the ones with almost no mispredicts to make cheaper.

\subsection{v2: two more items from the same budget}

The stall line in the budget was small on the kernels but suspicious. It was caused by the same
omission that produced the first defect above: the M$\rightarrow$X forwarding refused loads
because the forwarded value had no memory-data case. Supplying the memory output directly
removes the stall altogether. The hazard now fires only where the operand must be consumed in
decode, which is a branch or a \texttt{jalr} comparison.

\texttt{jalr} also moved into decode. Its target is \texttt{rs1} plus the immediate, both
available there, and because the new flush clears only the F/D register the \texttt{jalr}
continues into execute and still writes its link register. A resolved bit rides along to
execute so that the second check does not redirect it again, which would cancel the gain.

\begin{center}
\begin{tabular}{lrrrr}
\toprule
kernel & v1 & + load-use & + jalr & why \\
\midrule
\texttt{array-idx} & 88 & \textbf{80} & 80 & 8 stalls were 9\% of its cycles \\
\texttt{sum-array} & 48 & \textbf{44} & 44 & 4 stalls were 8\% \\
\texttt{callret}   & 34 & 34 & \textbf{33} & its one \texttt{jalr} miss \\
\textbf{total}     & \textbf{1288} & 1276 & \textbf{1275} & $-1.0\%$ against v1 \\
\bottomrule
\end{tabular}
\end{center}

The stall counter is zero on every test in the suite afterwards. Against the original baseline
this build is 1339 to 1275 cycles, $-4.8\%$.

\subsection{Graded result, and a six-fold surprise}

\begin{center}
\begin{tabular}{lrrrc}
\toprule
graded program & October & + prediction & + early resolution & + load-use and \texttt{jalr} \\
\midrule
Perf        & 0.587 & 0.499 & 0.487 & \textbf{0.453} \\
Fibonacci   & 0.586 & 0.499 & 0.489 & \textbf{0.474} \\
Bubble sort & 0.248 & 0.236 & 0.231 & \textbf{0.196} \\
Factorial   & 0.147 & 0.127 & 0.123 & \textbf{0.113} \\
GCD         & 0.084 & 0.069 & 0.068 & \textbf{0.065} \\
Sum array   & 0.057 & 0.054 & 0.052 & \textbf{0.045} \\
Check vowel & 0.183 & 0.174 & 0.173 & 0.175 \\
Swap        & 0.016 & 0.016 & 0.015 & 0.014 \\
Swap shift  & 0.017 & 0.016 & 0.016 & 0.015 \\
Simple if   & 0.013 & 0.012 & 0.012 & 0.012 \\
Simple add  & 0.012 & 0.011 & 0.011 & 0.011 \\
\midrule
\textbf{total (s)} & \textbf{1.950} & \textbf{1.713} & \textbf{1.677} & \textbf{1.573} \\
vs October  & -- & $-12.2\%$ & $-14.0\%$ & $\mathbf{-19.3\%}$ \\
\bottomrule
\end{tabular}
\end{center}

All eleven programs correct. Between the two early-resolution builds, eight are faster, two
unchanged and one is two milliseconds slower, which is the grader's quantum rather than a
regression.

The expensive part of v2 was worth 1.0\% on the project's own kernels and 6.2\% on the graded
programs. The mechanism is the same in both cases, one cycle per load-use pair, so the
difference is the workload: the hand-written kernels are arithmetic loops with a handful of
loads, while bubble sort and sum array spend their inner loops loading array elements. Those two
move by 15.2\% and 13.5\%.

\section{What the five experiments add up to}

\begin{center}
\small
\setlength{\tabcolsep}{5pt}
\begin{tabular}{llrr}
\toprule
change & mechanism & graded effect & local effect \\
\midrule
pre-decoder + BHT/BTB & remove mispredicts & $-12.2\%$ & about $-20\%$ \\
macro-op fusion v1    & remove instructions & $+0.4\%$ & $-21.7\%$ \\
macro-op fusion v2    & remove instructions & $+0.5\%$ & $-22.7\%$ \\
early resolution      & make mispredicts cheaper & $-2.1\%$ & $-3.8\%$ \\
load-use and \texttt{jalr} & remove stall cycles & $-6.2\%$ & $-1.0\%$ \\
dead-logic removal & smaller model, same cycles & $\mathbf{-1.7\%}$ & $0\%$ \\
\midrule
cumulative & & $\mathbf{-20.7\%}$ & \\
\bottomrule
\end{tabular}
\end{center}

Two features were measured and rejected, two measured and kept. The table also carries the most
transferable lesson of the project.

The local suite overstated fusion by roughly fifty times. It understated the load-use work by a
factor of 1.3. A hand-written kernel suite can be wrong in both directions, and the only way to
know which is to measure on the workload that is actually scored. That is not a reason to skip
the local suite: the sweep family, the fusion census and the 21.4\% bound are what made the
fusion experiment cheap to run and its failure interpretable. It is a reason to treat the
autograder as part of the experiment loop rather than as a final check.

\subsection{What the autograder actually measures}

One anomaly survived everything else in this report. Fusion removed 21.7\% of the cycles on the
kernels and changed the graded total by nothing. The obvious explanation was that the pattern
is absent from the compiled binaries, and the second variant was built to test it; when that
also returned nothing, the explanation was in doubt.

There is a second explanation, and it is testable without changing a single cycle. The graded
runtime is a simulation, so it is proportional to \emph{cycles} times \emph{the cost of
evaluating one cycle}, and the second term grows with the amount of logic in the model. Fusion
added a recogniser, a payload on two pipeline registers and a second combinational
instruction-memory read. If the second term is real, then deleting that logic while keeping the
cycle count identical should make the graded runtime fall.

The lean build is exactly that experiment: fusion removed, everything else untouched,
**1275 cycles either way**, all 69 local tests passing, all fourteen kernels at identical cycle
counts.

\begin{center}
\small
\begin{tabular}{lrrrr}
\toprule
graded program & October & v2 build & lean build & vs v2 \\
\midrule
Fibonacci   & 0.586 & 0.474 & \textbf{0.458} & $-3.4$\% \\
GCD         & 0.084 & 0.065 & \textbf{0.063} & $-3.1$\% \\
Check vowel & 0.183 & 0.175 & \textbf{0.170} & $-2.9$\% \\
Factorial   & 0.147 & 0.113 & \textbf{0.112} & $-0.9$\% \\
Bubble sort & 0.248 & 0.196 & \textbf{0.195} & $-0.5$\% \\
Perf        & 0.587 & 0.453 & \textbf{0.452} & $-0.2$\% \\
Sum array   & 0.057 & 0.045 & 0.045 & $0.0$\% \\
Simple if   & 0.013 & 0.012 & 0.011 & $-8.3$\% \\
Swap, swap shift, simple add & & & unchanged & $0.0$\% \\
\midrule
total (s)   & 1.950 & 1.573 & \textbf{1.546} & $-1.7$\% \\
\bottomrule
\end{tabular}
\end{center}

$27\,\mathrm{ms}$ on a build with identical cycles. All programs still correct.

So the metric has two terms and both are real. It also retro-explains the ledger: every feature
was a mix of cycles saved and model size added\begin{itemize}
\item prediction added tables and saved about 20\% of the cycles, so the cycle term dominated;
\item fusion added a recogniser and a second memory port to save 21.7\%, and the two terms
  cancelled exactly;
\item early resolution added a comparator to save 3.8\%, netting $-2.1$\%;
\item load-use forwarding added one multiplexer and saved the stall, netting $-6.2$\%;
\item deleting dead logic added nothing and saved nothing in cycles, netting $-1.7$\%.
\end{itemize}

The practical consequence is uncomfortable but useful: **a cycle saved and a gate added can
cancel on this metric**, so the cheapest way to score well is to make the design smaller at
equal cycles, not only faster. That is worth stating plainly in the report of a project whose
stated goal is a faster core, because the two objectives are not identical.

\section{Why the fractional part of CPI is a ceiling}

Everything in the previous section is bounded by one identity. With a single-issue pipeline
that retires at most one instruction per cycle,

$$\mathrm{CPI} = \frac{\text{cycles}}{\text{instructions}} = 1 + \frac{\text{overhead cycles}}{\text{instructions}} \;\ge\; 1$$

so the whole of the remaining headroom is the \emph{fractional part} of CPI, and every
optimisation in this report has been an attack on that fraction. CPI below 1 is not reachable
by making the fraction smaller; it needs more than one retirement per cycle, which is a
different machine.

\subsection{The fraction shrinks as programs get longer}

The overhead is not proportional to program length: it is mostly a fixed cost per run plus a
cost per loop, and the per-instruction share of a fixed cost falls as the program grows.

\begin{center}
\small
\begin{tabular}{lrrrr}
\toprule
kernel & instructions & cycles & CPI & overhead cycles \\
\midrule
\texttt{smoke}         &   2 &   6 & 3.00 &  4 \\
\texttt{sweep4}        &  19 &  26 & 1.37 &  7 \\
\texttt{callret}       &  25 &  33 & 1.32 &  8 \\
\texttt{loop10}        &  37 &  44 & 1.19 &  7 \\
\texttt{if-in-loop}    &  40 &  50 & 1.25 & 10 \\
\texttt{fib}           &  59 &  66 & 1.12 &  7 \\
\texttt{array-idx}     &  73 &  80 & 1.10 &  7 \\
\texttt{nested}        &  88 & 102 & 1.16 & 14 \\
\texttt{array-swap-loop} & 104 & 129 & 1.24 & 25 \\
\texttt{sweep50}       & 158 & 164 & 1.04 &  6 \\
\texttt{calls-in-loop} & 159 & 196 & 1.23 & 37 \\
\texttt{smul}          & 205 & 213 & 1.04 &  8 \\
\texttt{sweep100}      & 308 & 314 & 1.02 &  6 \\
\bottomrule
\end{tabular}
\end{center}

Read the last column, not the CPI. \texttt{sweep100} executes 308 instructions and wastes six
cycles in total: four of pipeline fill and two redirects. \texttt{sweep4} executes 19
instructions and wastes seven, for a CPI of 1.37. The waste is almost the same in both; only the
denominator changed.

That is the ceiling. Even a machine with no stalls and no mispredicts left to remove could take
\texttt{sweep100} from 314 cycles to 308, which is 2\%. \texttt{sweep4} would improve 27\%, but
\texttt{sweep4} contributes 26 cycles out of 1600 to the suite's total, so it barely matters.

\subsection{What the six cycles actually are}

For \texttt{sweep100}, six cycles of overhead across 308 instructions:

\begin{itemize}
\item \textbf{four of pipeline fill}. The first instruction takes four cycles to reach writeback
  and those cycles are counted. A five-stage pipeline cannot do better without being a
  four-stage pipeline.
\item \textbf{two of redirect}, one for the cold first iteration and one for the loop exit. The
  entry miss is unavoidable: nothing can predict a branch before it has been seen. The exit
  miss is what a loop-exit predictor would remove, worth one cycle per loop.
\end{itemize}

The larger fractions in the table are the kernels with structure, not the kernels that are long:
\texttt{calls-in-loop} wastes 37 cycles because twenty returns through two call sites alternate
between targets, and \texttt{array-swap-loop} wastes 25 because a branch waits for a load. Both
are the structural costs covered earlier, and both are per-iteration, not fixed, which is why
they survive into long programs where the fixed costs have vanished.

\subsection{Amdahl applied to the fraction}

The total across the sixteen kernels is 1600 cycles for 1433 instructions, a CPI of 1.12, so
the fraction is 10.4\%. An optimisation that removed \emph{every} wasted cycle would therefore
give at most 10.4\%, and most of the fraction is not removable: fill is structural, and the
entry mispredict cannot exist without a predictor that has seen the branch before.

So the realistic statement is: the suite-level ceiling is around 10\%, the addressable part is
perhaps half of it, and it shrinks as programs get longer, which the graded programs are.

\subsection{There is a second ceiling, in the measurement}

The autograder reports milliseconds on programs of 11\,ms to 500\,ms, so its resolution is 8\%
on the smallest and 0.2\% on the largest. Repeating the same build produces differences of
1--3\,ms on the small programs, which is up to 3\%. The consequence is blunt: \textbf{graded
differences below about 1\% cannot be measured at all}, which is exactly what the last two
experiments returned --- the memory X-checks moved the total by 4\,ms, with two programs faster
and two slower.

Every real result in this report is comfortably above that floor: prediction 12.2\%, load-use
6.2\%, early resolution 2.1\%, dead logic 1.7\%. Anything smaller would have been
indistinguishable from noise, and treating it as a result would have been a mistake.

\subsection{What would break the ceiling}

\begin{center}
\begin{tabular}{lll}
\toprule
change & effect on the identity & cost \\
\midrule
shorter pipeline & fill falls from 4 to 2--3 & a redesign of forwarding and hazards \\
dual issue & $\mathrm{CPI}$ can fall below 1 & months; doubles most of the datapath \\
loop-exit predictor & removes 1 cycle per loop & days; helps only many-short-loop code \\
return-address stack & removes return mispredicts & days; a teammate's branch \\
\bottomrule
\end{tabular}
\end{center}

The first two change the identity itself. The last two only shrink the fraction, and the table
above says how much room is left there: on the long kernels, almost none.

\section{The synthetic suite against the graded runs}

Every feature in this report was measured twice: on sixteen kernels written for this project, and
on the eleven programs the autograder runs. The two measurements were not equally useful, and the
pattern in their disagreement is the most transferable result in the document.

\subsection{The two measurements}

\begin{center}
\small
\begin{tabular}{lll}
\toprule
& synthetic & graded \\
\midrule
what runs & 16 hand-written kernels & 11 compiled programs, sources unknown \\
what is reported & cycles, stalls, redirects, CPI & wall time of the simulation, rounded to 1\,ms \\
resolution & one cycle, always & 1\,ms; 8\% on the smallest program \\
what it isolates & machine behaviour & machine behaviour $\times$ simulator cost \\
\bottomrule
\end{tabular}
\end{center}

They are not the same measurement, and neither substitutes for the other.

\subsection{Every feature, side by side}

\begin{center}
\footnotesize
\setlength{\tabcolsep}{4pt}
\begin{tabular}{lrrrr}
\toprule
feature & predicted & synthetic & graded & transfer \\
\midrule
pre-decoder                    & $\approx$$-4$ cycles/test & $-4$ to $-8$ & n/a (in all builds) & -- \\
BHT+BTB                        & $-2$ cycles/mispredict & $-20.3\%$ on \texttt{loop10} & $-12.2\%$ & 0.6 \\
fusion v1                      & $-21.4\%$ bound & $-21.7\%$ & $+0.4\%$ & $\approx 0$ \\
fusion v2 (both operands)      & $>-21.4\%$ & $-22.7\%$ & $+0.5\%$ & $\approx 0$ \\
early resolution               & $\approx -4\%$ & $-3.8\%$ & $-2.1\%$ & 0.6 \\
load-use removal + \texttt{jalr} & $\approx -1\%$ & $-1.0\%$ & $-6.2\%$ & 6.2 \\
dead-logic removal             & $0$ cycles & $0$ & $-1.7\%$ & $\infty$ \\
memory X-checks                & $\approx -1\%$ & $0$ & $-0.3\%$ & noise \\
\bottomrule
\end{tabular}
\end{center}

The transfer column is graded effect divided by synthetic effect. It ranges from zero to
infinity, which is the point: a synthetic measurement tells you what the machine did, not what
the grade will be.

\subsection{Three reasons they diverge, in order of how much they cost}

\paragraph{The workload is different.} The kernels are hand-written loops with a branch in the
body. The graded programs are compiled C and spend their inner loops loading array elements and
calling subroutines. Fusion is exactly the optimisation that a hand-written loop overstates: a
counted loop written in the addi-then-branch form presents one fusable pair per iteration,
guaranteed, while compiled code rarely presents that adjacency at all. Load-use is the opposite
case, and it is why the same pair of numbers reads $-1.0\%$ synthetic against $-6.2\%$ graded.

\paragraph{The metric has a second term.} The graded runtime is a simulation, so it is
proportional to cycles \emph{and} to the cost of evaluating a cycle, which grows with the logic
in the model. Fusion removed 21.7\% of the synthetic cycles and added a recogniser, a payload on
two registers and a second memory read; the graded total did not move. The lean rebuild is the
control experiment: same cycles, less logic, $-1.7\%$.

\paragraph{The measurement is coarser and noisier.} One millisecond on an 11\,ms program is 8\%.
Repeated runs of an unchanged build differ by 1--3\,ms. The last experiment produced a 4\,ms
change with two programs faster and two slower, which is noise, and is recorded as noise.

\subsection{What the synthetic suite is still for}

It caught every regression, and it is the only place where a change can be attributed at all:
one cycle is one cycle, with no rounding and no noise. The three bring-up bugs in the early
resolution work were each found by a synthetic trace, not by a graded run. It also produced the
21.4\% bound that made the fusion experiment cheap to run and its failure interpretable ---
without that bound, a graded result of $+0.4\%$ would have looked like a broken build.

What it cannot do is rank features against each other, because the ranking depends on a workload
it does not contain. That is a structural limitation, not a bug in the kernels, and the only fix
is to measure on the graded programs often enough that the disagreement shows up early.

\section{Engineering notes}

A few things cost more time than the RTL did, and are worth writing down.

\paragraph{A regression that was correct but saved nothing.} Leaving the branch opcode in the
execute-stage prediction check makes both mechanisms redirect the same branch, with execute
winning. Every test passed. The cycle count was identical to the baseline. Any check that only
looks at pass/fail would have concluded the feature did not help, when in fact the feature was
not running.

\paragraph{Three false greens.} A unit testbench that counted its rows in a second
\texttt{initial} block could sample zero rows and report success; the fix is to count at the end
of the block that fills the table, with a guard that aborts on an empty table. A test that
placed its failure label immediately after \texttt{ecall} had it execute during the five-cycle
drain window and overwrite the pass flag, so a correct design reported failure. And verifying a
submission zip by extracting it and running the suite silently reused the build artefacts that
came with it, testing the old binary; the fix is to delete the build directory first.

\paragraph{Counters that lie.} The squashed-instruction counter falls from 52 to 1 when branches
resolve in decode. That is not a 50-fold improvement. The counter appears to count execute-stage
flushes, and a decode-stage flush clears one register. Cycles are the metric; the counter is
context.

\paragraph{The wrong debugging tool.} A differential PC trace found the fusion integration bug
by showing where two builds diverged. The same method is useless for early resolution, because
that change deliberately reduces how many instructions are fetched speculatively, so the PC
sequences differ for correct reasons. Comparing architectural writebacks instead finds the first
wrong value directly and found all three bring-up defects.

\paragraph{Two mechanisms, one redirect.} Both redirect paths must be gated by the same feature
mask, or one quietly overrides the other. This is worth a comment in the RTL, and it is why the
execute-stage check now names the reason it excludes branches.

\section{Current state}

The suite is 67 tests plus one long-standing failure in \texttt{sample-bmarks/default}, which has
failed since 2026-09-23 and is unaffected by any of this work. All 67 pass on the shipped
configuration.

The shipped configuration is: pre-decoder, BHT with tags and a valid bit, BTB, early branch
resolution, load-use forwarding and \texttt{jalr} resolution in decode. The fusion path has been
removed from the design outright: the recogniser, its payload ports on both pipeline registers,
the fused comparator in execute and the fetch stride input are all gone. The module and its
testbench remain in the tree, the module marked as not instantiated, so the measurement stays
reproducible.

Not measured: area and critical path. Synthesis requires the department machines. The synthesis
script has been changed to blackbox the instruction memory, because elaborating the array turns
it into tens of thousands of flip-flops whose read path dominates the critical path, which is
what an earlier attempt measured. A clock-period sweep script is in place so that the maximum
frequency can be bracketed in one run rather than one run per constraint. Until that is done,
the time figures in this report are cycle counts, and the wall-clock gains assume the clock did
not get slower.

The next candidate for the remaining mispredict cost is a return-address stack. The
\texttt{callret} benchmark exists for exactly that measurement and has never been used as one.

\appendix
\section{Benchmark kernels}

\begin{center}
\small
\begin{tabular}{lll}
\toprule
kernel & shape & self-check \\
\midrule
\texttt{smoke}     & two instructions & flow only \\
\texttt{sweep4/10/25/50/100} & counted loop, one branch, five trip counts & iterations $= N$ \\
\texttt{loop10}    & counted loop with accumulation & $1+\cdots+10 = 55$ \\
\texttt{nested}    & 5 outer, 4 inner & sum $= 30$ \\
\texttt{if-in-loop}& hot loop, one fall-through & marker $= 1$ \\
\texttt{sum-array} & store/load loop over four words & total $= 10$ \\
\texttt{array-idx} & \texttt{s += a[i]} with a scaled index & total $= 36$ \\
\texttt{fib}       & iterative Fibonacci & $\mathrm{fib}(11) = 89$ \\
\texttt{smul}      & shift-and-add multiply, 32 iterations & $6 \times 7 = 42$ \\
\texttt{callret}   & three calls and three returns & $3 \times 5 = 15$ \\
\bottomrule
\end{tabular}
\end{center}

\section{Reproducing}

\begin{verbatim}
    make sim-all -C verif/               # 67 tests, every directory
    make bench-build -C verif/           # assemble every benchmark source
    make tb-pre_decode_tb -C verif/      # unit testbench, 27 rows
    make tb-fuse_detect_tb -C verif/     # unit testbench, 16 rows
    python3 code/benchmarks/fusion_analysis.py
    python3 code/benchmarks/objdump_analysis.py <listing.objdump>
\end{verbatim}

Every measurement in this report is a full \texttt{make sim-all} on the branch named in the
corresponding experiment section, with the feature mask stated. The submission zips named in
sections 4 and 5 are kept alongside the sources.

\end{document}
